% intro_paragraphs.tex -- SOTA literature-review paragraphs for the OREM
% manuscript Introduction (issue #46, broader background survey).
% Companion files: literature_matrix.md, references.bib (this folder).
%
% This is the broad state-of-the-art framing paragraph(s); the specific
% novelty paragraph distinguishing this work from Mutyalarao & Sharma
% (2010,2011) / Sellamuthu & Sharma (2017,2018) / Sharma & Raj is already
% drafted separately in issue #46's own comment thread and should follow
% immediately after this block in the final Introduction, not duplicate it.

The prediction of atmospheric re-entry for resident space objects (RSOs)
decaying from highly eccentric orbits (HEO, $e>0.2$) and geostationary
transfer orbits (GTO) is a classical problem in astrodynamics and Space
Domain Awareness (SDA), originating with Shute's early treatment of
highly-eccentric-satellite lifetimes \cite{shute1965,shutechiville1966}
and King-Hele's foundational analytical theory of orbital contraction
under upper-atmospheric drag dissipation \cite{kinghele1987}. Building on
this lineage, Sharma and coworkers developed a family of
Kustaanheimo--Stiefel (KS) regularized analytical and semi-analytical
theories that couple drag-driven semi-major-axis reduction with Earth's
oblateness \cite{sharmaraj1988,sharma1997} and the diurnal atmospheric
density bulge \cite{swinerdboulton1983}, establishing regularized orbital
elements as a numerically robust alternative to Cartesian propagation in
the highly eccentric, drag-active regime addressed by this work. More
recent, standards-level assessments confirm that atmospheric-density-model
uncertainty -- not propagator or force-model fidelity -- remains the
dominant error source for orbit-lifetime and re-entry-time prediction
across the field \cite{vallado2014,iso27852}.

Beyond drag, lunisolar third-body gravitational perturbations exert a
first-order influence on the long-term evolution of GTO/HEO perigee
altitude. Kozai's secular perturbation theory \cite{kozai1962} and Hough's
critical-inclination analysis \cite{hough1981} established that third-body
forcing can drive large-amplitude eccentricity and perigee-lowering
oscillations near the classical critical inclination, a mechanism
independently corroborated by Colombo's long-term, drag-free evolution
study of highly elliptical orbits \cite{colombo2015}. For GTO specifically,
Wang and Gurfil identified and characterized \emph{solar apsidal
resonance} -- a 1:1 commensurability between the GTO apsidal-line
precession rate and the Sun's apparent motion -- that converts the
normally periodic, few-hundred-day perigee-height oscillation into a
monotonic, one-way collapse or rise
\cite{wanggurfil2016,wanggurfil2017asr,wanggurfil2017jsr}, with Luo and
Wang subsequently generalizing this lunisolar resonance corridor to
inclined GTOs \cite{luowang2019}. These resonance corridors impose a
fundamental, quantifiable ceiling on deterministic re-entry predictability:
post-resonance orbital evolution becomes acutely sensitive to sub-degree
uncertainties in right ascension of the ascending node and area-to-mass
ratio, a sensitivity independently reproduced within a KS-regularized
dynamical framework closely related to the propagator used in this work
\cite{sellamuthusharma2021}.

The response-surface-methodology (RSM) and genetic-algorithm (GA) approach
used here to fit an object's eccentricity and ballistic number to its
observed apogee decay traces to Sharma, Bandyopadhyay, and Adimurthy's
original zone-fitting cost function \cite{sharmabandyopadhyay2006} and its
extension to short-term (${<}\,3$-month) GTO decay-zone prediction by
Mutyalarao and Sharma \cite{mutyalaraosharma2010,mutyalaraosharma2011}; a
structurally similar essential-ballistic-parameter fit was independently
developed by Gupta and Anilkumar for terminal-phase decay
\cite{guptaanilkumar2015}. A recent benchmarking study within this same
methodological lineage found that no single re-entry-prediction method --
RSM-GA included -- dominates across the full decay horizon, motivating
multi-method or ensemble approaches \cite{dutt2023}. This growing body of
SDA modeling work is set against a tightening regulatory landscape: the
Inter-Agency Space Debris Coordination Committee's mitigation guidelines
\cite{iadc2021} and, more recently, the U.S. Federal Communications
Commission's rule requiring post-mission disposal within five years of
end-of-mission for satellites passing through low-Earth orbit
\cite{fcc2022,fedreg2024} have sharply shortened the operationally relevant
re-entry-prediction horizon relative to the legacy 25-year disposal
guideline, increasing the practical value of accurate short- and
medium-term GTO/HEO decay prediction of exactly the kind addressed in this
work.
